%% ==========================================================================
%%  SECTION 4 --- PRIVATE TRAINING                              [Track T2]
%%
%%  Job: trace the private matrix-factorisation line to its current state of
%%  the art, and land ONE argument -- that the choice of ALGORITHM is a
%%  CRYPTOGRAPHIC decision, not a machine-learning one. ~2.5 pages.
%% ==========================================================================
\section{Private Training}
\label{sec:training}

\todo{Frame: everyone in this line computes the same thing --- a low-rank
      factorisation $\ratings \approx \userembed \cdot \itemembed$ --- and the
      entire history is about which ALGORITHM for computing it survives being
      run under cryptography. That is the section's thesis; state it first.}

\subsection{Matrix factorisation, in the clear}
\label{sec:training:mf}

\todo{Half a page, no more. $\ratings \in \R^{\users \times \items}$,
      embedding dimension $\dime \ll \min(\users,\items)$, minimise
      $\norm{\ratings - \userembed\cdot\itemembed}_F$. Note that when
      $\itemembed$ is orthonormal, $\userembed = \ratings \cdot \itemembed^\top$,
      so recommendations for user $i$ depend only on $\uvec{i}$ and
      $\itemembed$. \verify{\nudgeplain{} \S2} That fact does a lot of work
      later --- flag it here and return to it in \S8.}

\subsection{Garbled circuits: Nikolaenko et al.}
\label{sec:training:gc}

\todo{CCS'13. Gradient descent inside a garbled circuit. Establish it as the
      baseline the field has been escaping ever since --- \nudgeplain{} reports
      beating this family by FOUR ORDERS OF MAGNITUDE in communication and
      compute. \verify{exact claim}. Say why: data-oblivious gradient descent
      forces a circuit whose size scales with the whole rating matrix.}

\subsection{\pirsonaplain{}: four-party Boolean factorisation}
\label{sec:training:pirsona}

\todo{PoPETs'21. \pirsona{} --- the instructor's paper, so be thorough.
      Cover: the 4PC Boolean matrix factorisation; the new primitives it
      contributes (4PC fixed-selection MUX/DeMUX, fast 3PC vector
      normalisation for shared fixed-point vectors, 4PC well-formedness tests
      for $(2,2)$-DPFs, one-round 3PC integer comparison); and the pairwise
      non-collusion assumption it pays for them.
      Its DELIVERY half belongs in \S6, not here --- cross-reference, do not
      duplicate.}

\subsection{\nudgeplain{}: power iteration as a matrix-vector program}
\label{sec:training:nudge}

\todo{USENIX Sec'26. The current state of the art, and the section's
      centrepiece. Build the argument in this order:
      (1) recall \Cref{obs:prim:matvec-free} --- matrix-vector work is free;
      (2) therefore prefer an algorithm that is mostly matrix-vector products
          separated by cheap non-linear steps;
      (3) power iteration has exactly that shape and gradient descent does
          not. THIS IS THE POINT OF THE WHOLE SECTION.
      Then the matrix-vector program abstraction (Def 4.1) and Thm 4.2.}

\begin{algorithm}[t]
\caption{$\ApproxFactor$: matrix factorisation by power iteration \cite{nudge2026}}
\label{alg:training:approxfactor}
\todo{Transcribe \nudgeplain{} Fig.\ 4 faithfully, with a citation. Mark clearly
      which lines are FREE (the two $\mathrm{Mul}$ calls) and which are
      INTERACTIVE ($\Normalize$). Annotate the line where each row of
      $\itemembed$ is OPENED IN THE CLEAR --- that is the design's load-bearing
      trick and the hinge of \S8.}
\end{algorithm}

\subsubsection{The two protocols that cost anything}
\label{sec:training:nonlinear}

\todo{$\Trunc_{\fracbits}$: deterministic truncation with sign extension;
      low-order carries fixed by an integer comparison, high-order carries
      prevented by one bit of slack. 3 rounds,
      $2\dime\fracbits(\seclen+4)+10\dime\bits$ bits. Emphasise WHY that is an
      improvement: the leading term is $\seclen\fracbits$ rather than
      $\seclen\bits$, giving $\approx\bits/\fracbits = 6\times$ at
      $\bits=\seclen=128$, $\fracbits=20$.
      $\ApproxNormalize$: shares of $1/\norm{v}$. $\norm{v}^2$ is degree-two
      hence free; the inverse square root is seeded from the
      most-significant-non-zero-bit via $\bits+1$ SIMULTANEOUS comparisons in
      one round, then Newton--Raphson. Explain what this avoids --- the
      $O(\bits)$-round route and the exponential-lookup-table route.
      \verify{all constants against \nudgeplain{} \S4.3, Lemmas 4.3 and 4.4}}

\subsection{The non-cryptographic alternatives}
\label{sec:training:alternatives}

\todo{Short --- one paragraph each, for contrast and completeness.
      Federated MF (Chai et al.) and its leakage through gradients;
      FHE-based MF (arXiv 2509.03024) and its cost;
      Paillier/decentralised variants (PADER, 2601.10212);
      recent secret-sharing recommenders (PRSI, 2411.18653).
      The purpose here is to show we surveyed the field, and to justify why
      \S8 targets the MPC line.}

\begin{table}[t]
\centering
\todo{Table: system $\times$ \{parties, corruption model, algorithm, dataset,
      scale reached, cost, quality metric\}. Include Nikolaenko,
      \pirsonaplain{}, \nudgeplain{}, federated, FHE. Numbers must carry their
      SETTING --- hardware and network --- or they are not numbers.}
\caption{Private training systems compared.}
\label{tab:training:compare}
\end{table}

\subsection{Where training now stands}
\label{sec:training:standing}

\todo{Close with: private training is essentially SOLVED at real scale ---
      Netflix, 0.5M users, 50 minutes, quality on par with non-private MF.
      Then the pivot that sets up \S5, quoting \nudgeplain{} \S3.1 verbatim:
      it ``relies on other means\dots{} to let users fetch data items in a
      private way.'' End the section on that quote.}
